\chapter{Discusión}\label{sec:discusion}

La comparación muestra que concatenar descriptores puede mejorar el desempeño respecto de elegir un único bloque, pero el beneficio depende de la composición y del dominio. En DTD y FMD, seleccionar un subconjunto ofrece resultados competitivos frente a utilizar toda la biblioteca. En CUReT, Completa alcanza los mayores valores; en Outex, la mejor estrategia depende del clasificador: Completa con SVM y GFS con ResMLP. La concatenación completa mejora descriptivamente frente a Individual en tres de los cuatro datasets al promediar los clasificadores, pero pierde desempeño en FMD. Seleccionar bloques ofrece una alternativa de menor dimensión, cuyo resultado depende de su composición. Este patrón es compatible con distintas combinaciones de información complementaria y redundante; el diseño no mide directamente esa redundancia ni identifica un mecanismo causal.

Las pruebas de Friedman, Quade y Holm moderan la interpretación de las ventajas descriptivas de la biblioteca de 21 descriptores. Con cuatro datasets no se obtuvo evidencia suficiente para afirmar una superioridad global. Esto no demuestra igualdad de rendimiento; delimita lo que sostienen los contrastes realizados.

\section{Limitaciones}

GFS es una búsqueda local y no garantiza el subconjunto óptimo. Top-$k$ hereda su presupuesto de bloques, por lo que compara composiciones bajo un $k$ compartido. Individual se elige internamente y no representa al descriptor con mejor resultado a posteriori en el test. Las particiones de DTD y FMD reutilizan imágenes, CUReT presenta dos direcciones complementarias y Outex tiene una sola partición. El número de datasets limita la precisión y la generalización de la inferencia multibase. No se realizó un cálculo de tamaño muestral a priori ni se presentan intervalos de confianza de generalización a nuevos datasets. Las desviaciones entre particiones describen la variabilidad observada; no estiman esa incertidumbre tratando folds dependientes como réplicas independientes.

Los descriptores clásicos se calcularon sobre escala de grises, mientras las redes y RGB Pixel N-grams utilizan RGB; en CUReT las imágenes de origen son grises. Además, los corpus de preentrenamiento de los modelos podrían solaparse con imágenes de los benchmarks; ese solapamiento no fue auditado. El agrupamiento de conteos y la reducción de dimensión pueden perder información. El histograma restringido de LBP-rot también descarta códigos; la comparación incluye esa implementación específica y no permite extraer conclusiones sobre el DRLBP original. Los resultados dependen de la parametrización especificada. No se dispone de mediciones homogéneas de latencia y memoria para atribuir ahorro de cómputo a la reducción dimensional.
